\section{Amplifying the collar estimate}\label{sec:amplification}

Proposition~\ref{prop:collar} gives a sublinear bound on the mutual
information between a target and the region beyond its collar.  We now
make this mutual information smaller than any fixed inverse power of
the scale, at the cost of a slightly wider collar.  Typical spectral
subspaces first give a candidate restoring operator whose norm is
controlled by the collar mutual information.  Positive constraints
suppress its excited component, yielding accurate restoration after
localization.  A polar decomposition then replaces this operator by a
local unitary.  An untouched reference copy records the target's
correlations with the remote system and witnesses their decoupling.

Throughout this section fix
\begin{equation}\label{eq:amplification-exponents}
 \varepsilon=10^{-5},\qquad \beta=1-2\cdot10^{-6},\qquad
 \alpha=\frac{p}{p+1}>1-10^{-6},
\end{equation}
where the integer $p$ is fixed.  Apply Proposition~\ref{prop:positive}
with this value of $p$.  Thus, in this section, $H_F=\sum_i k_i$ is a
fresh exact positive system with
\begin{equation}\label{eq:amplification-positive-system}
 0\le k_i\le I,\quad k_i\Omega=0,\quad
 H_F\ge g(I-P_\Omega),\qquad
 P_\Omega=|\Omega\rangle\langle\Omega|,
\end{equation}
and $g>0$ is independent of the finite domain.  The scanner's
truncated Hamiltonian is no longer used.  All neighborhoods below are
neighborhoods in the induced graph: for $S\subseteq\Lambda$ we write
$N_r(S)=\{y\in\Lambda:d_\Lambda(y,S)\le r\}=N_r^\Lambda(S)$.

\subsection{Localizing products of positive constraints}

The following estimate supplies the localization needed later.  It is
stated for a general operator, which may act on additional finite
systems.  These additional systems are spectators in the dynamics.
The spatial truncation argument has the propagation structure of
dissipative Lieb--Robinson estimates~\cite[Theorems~1 and~2]{BarthelKliesch2012}.
We prove it here on coupled clock realizations, since the needed bound
is an expectation of a norm.

\begin{lemma}[Local constraint products]\label{lem:amplification-local-products}
Let positive contractions $k_i$, anchored at $a_i\in\Lambda$ with
bounded multiplicity, satisfy
\[
 \|k_i-E_{i,l}k_i\|\le C e^{-c l^\alpha}\qquad(l\ge0),
\]
where $E_{i,l}$ is conditional expectation onto $N_l(a_i)$.
Set $G_i=(I-k_i)^{1/2}$ and $K_i=k_i^{1/2}$.
Fix $t\ge0$ and an integer $r\ge1$.
At each anchor label run an independent rate-one Poisson clock on
$[0,t]$.  In occurrence order, with later factors on the left, let
$D$ be the product of all $G_i$, and let $D_{\rm ret}$ retain only
events with $d_\Lambda(a_i,S)\le r$, for a nonempty set $S\subseteq\Lambda$.
Let $L$ be the same retained product with each $G_i$ replaced by
$G_{i,r}=(I-E_{i,r}k_i)^{1/2}$.

There are uniform constants $C,c,v>0$ such that
\begin{equation}\label{eq:amplification-product-truncation}
 \mathbb E\|D_{\rm ret}-L\|
 \le Ct|S|(1+r)^2e^{-c r^\alpha}.
\end{equation}
In particular $L$ is a contraction supported on $N_{2r}(S)$.
If $k_i\Omega=0$ and an operator $B$ is supported on $S$ and
spectator systems, then, for every unit spectator vector $\xi$,
\begin{equation}\label{eq:amplification-product-omission}
 \mathbb E\|(D-D_{\rm ret})B(\Omega\otimes\xi)\|
 \le Ct|S|\|B\|e^{vt-c r^\alpha}.
\end{equation}
Both bounds concern the same coupled clock realization.  Their constants
are independent of $|\Lambda|$ and of spectator dimensions.
\end{lemma}

\begin{proof}
We first prove an operator estimate for unital channels.  Define
\[
 \mathcal E_i(B)=G_iBG_i+K_iBK_i,
 \qquad
 \mathcal E_{i,l}(B)=G_{i,l}BG_{i,l}+K_{i,l}BK_{i,l},
\]
where $K_{i,l}=(E_{i,l}k_i)^{1/2}$.  When $k_i\Omega=0$, functional
calculus gives $G_i\Omega=\Omega$ and $K_i\Omega=0$, so
\begin{equation}\label{eq:amplification-channel-vector}
 \mathcal E_i(C)(\Omega\otimes\xi)=G_iC(\Omega\otimes\xi)
\end{equation}
for every operator $C$ and spectator vector $\xi$.  Products of these
channels therefore recover the contraction products when evaluated on
the ground vector.  We localize the channels first: their unitality
makes each local channel fix the algebra outside its supporting ball.
They are also completely positive, hence contractions in operator norm,
including after adjoining spectators.  Lemma~\ref{lem:quasilocal-roots} gives
\begin{equation}\label{eq:amplification-channel-approximation}
 \|\mathcal E_i-\mathcal E_{i,l}\|_{\rm cb}
 \le C e^{-c l^\alpha}.
\end{equation}
The completely bounded norm here is the supremum of the operator norms
after adjoining identity maps on finite spectator systems.  In this
case the estimate follows directly from the norm bounds for the two
square roots, so it has no dimension factor.

Put $\mathcal E_{i,-1}=\id$ and
$\Delta_{i,l}=\mathcal E_{i,l}-\mathcal E_{i,l-1}$ for $l\ge0$.
By increasing $C$ to include $l=0$,
\begin{equation}\label{eq:amplification-channel-shells}
 \|\Delta_{i,l}\|_{\rm cb}\le C e^{-c l^\alpha}.
\end{equation}
Both channels in this difference fix operators commuting with the
algebra on $N_l(a_i)$.  Thus the difference vanishes on that
commutant.  For an operator $B$ define its oscillation at a physical
site $y$ by
\[
 d_y(B)=\sup_{U_y\ {\rm unitary}}\|[U_y,B]\|.
\]
Average over the on-site unitaries in a set $Q$, one site at a time.
The total norm change is at most $\sum_{z\in Q}d_z(B)$, and the result
commutes with the algebra on $Q$.  Consequently
\begin{equation}\label{eq:amplification-shell-oscillation}
 \|\Delta_{i,l}(B)\|
 \le C e^{-c l^\alpha}\sum_{z\in N_l(a_i)}d_z(B).
\end{equation}
This argument also applies to operators acting on spectators.

To bound the oscillation after one event, suppose first that
$m=d_\Lambda(a_i,y)<\infty$.  The channel $\mathcal E_{i,m-1}$
does not act at $y$ and commutes with conjugation by $U_y$; hence it
cannot increase $d_y$.  Telescope the shells from radius $m$ onward
and use $d_y(C)\le2\|C\|$ to obtain
\begin{equation}\label{eq:amplification-oscillation-increment}
 d_y(\mathcal E_i B)\le d_y(B)+
 C\sum_{l\ge d_\Lambda(a_i,y)}e^{-c l^\alpha}
       \sum_{z\in N_l(a_i)}d_z(B).
\end{equation}
For $m=0$ the initial channel is the identity.  If $m=\infty$,
$\mathcal E_i$ acts in a different component and cannot increase
$d_y$.  Exact component support follows either from the construction
or by letting $l$ exceed the component diameter in the approximation
assumption.

Sum the increment in
\eqref{eq:amplification-oscillation-increment} over the retained
anchor labels.  Define the nonnegative matrix
\[
 A_{yz}=C\sum_{i\ {\rm retained}}\sum_{l\ge0}
 e^{-cl^\alpha}\mathbf1_{\{y,z\in N_l(a_i)\}},
\]
with the constant from \eqref{eq:amplification-oscillation-increment}.
It has a uniformly bounded weighted row sum:
\begin{equation}\label{eq:amplification-weighted-kernel}
 \sup_y\sum_{z:\,d_\Lambda(y,z)<\infty}
 A_{yz}e^{a d_\Lambda(y,z)^\alpha}\le v
\end{equation}
for some $a>0$.  Indeed, at radius $l$ there are at most
$C(1+l)^2$ anchor labels whose balls contain $y$, and at most
$C(1+l)^2$ sites in each ball.  Their distances from $y$ are at most
$2l$.  The left side is therefore bounded by
\[
 C\sum_{l\ge0}(1+l)^4e^{-c l^\alpha+a(2l)^\alpha},
\]
which is finite for sufficiently small $a$.  There are no kernel
entries between different components.

Let $B_u$ be the result of the retained channel events up to time $u$,
starting from $B$.  The Poisson jump formula and
\eqref{eq:amplification-oscillation-increment} imply, in integral
form, the differential upper bound
\[
 \frac{d}{du}\mathbb E d_y(B_u)
 \le\sum_z A_{yz}\mathbb E d_z(B_u).
\]
The site and anchor sums are finite for a fixed domain, and the radius
sums converge absolutely with the uniform bounds in
\eqref{eq:amplification-weighted-kernel}.
Initially $d_y(B)\le2\|B\|1_S(y)$.  On components meeting $S$,
weight the last inequality by $e^{a d_\Lambda(y,S)^\alpha}$.
The triangle inequality and $0<\alpha<1$ give
\[
 d_\Lambda(y,S)^\alpha-d_\Lambda(z,S)^\alpha
 \le d_\Lambda(y,z)^\alpha.
\]
The weighted supremum thus grows at rate at most $v$.
Gronwall's inequality yields
\begin{equation}\label{eq:amplification-propagation}
 \mathbb E d_y(B_u)
 \le2\|B\|e^{vu-a d_\Lambda(y,S)^\alpha}.
\end{equation}
Oscillations on components disjoint from $S$ remain zero.

We can now remove the distant events.  For an omitted anchor at finite
distance $d=d_\Lambda(a_i,S)>r$, sum
\eqref{eq:amplification-shell-oscillation} over all $l\ge0$ and
apply \eqref{eq:amplification-propagation}.  When $l\le d/2$,
every $z\in N_l(a_i)$ has $d_\Lambda(z,S)\ge d/2$.
The sum over these radii is bounded by
$C\|B\|e^{vu-a(d/2)^\alpha}$.  For $l>d/2$ use the shell
decay and the ball cardinality instead.  Absorbing its polynomial
factor gives, in both cases,
\begin{equation}\label{eq:amplification-omitted-event}
 \mathbb E\|(\mathcal E_i-\id)(B_u)\|
 \le C\|B\|e^{vu-c'd^\alpha}.
\end{equation}
At infinite distance the left side is zero.

Couple full and retained processes by using the same clocks.
The difference of their products telescopes with earlier retained
channels and later full channels.  All the later channels contract
the operator norm.  The earlier retained process is independent of
the omitted clocks, so their expected contributions are obtained by
integrating the rate-one bound
\eqref{eq:amplification-omitted-event}.  There are at most
$C|S|(d+1)^2$ anchor labels at distance $d$ from $S$.
Summing the stretched-exponential tail and integrating time therefore
prove the coupled expected-norm estimate
\begin{equation}\label{eq:amplification-channel-omission}
 \mathbb E\|\mathcal E_{\rm full}(B)-\mathcal E_{\rm ret}(B)\|
 \le Ct|S|\|B\|e^{vt-c''r^\alpha}.
\end{equation}
In particular this is a bound on the expectation of the norm, not
merely on the norm of a difference of averaged operators.

If $k_i\Omega=0$, evaluating both channel products on
$\Omega\otimes\xi$ and iterating
\eqref{eq:amplification-channel-vector} converts
\eqref{eq:amplification-channel-omission} into
\eqref{eq:amplification-product-omission}.  Unitality is essential
in the channel argument: the single map $C\mapsto G_iCG_i$ need
not fix operators on a disjoint system.

Finally there are at most $C|S|(1+r)^2$ retained labels, and hence
their expected number of events is at most $Ct|S|(1+r)^2$.
Each event's contraction changes by at most $Ce^{-c r^\alpha}$
under radius-$r$ approximation.  Telescoping the contraction products
proves \eqref{eq:amplification-product-truncation}.  A retained
anchor is within $r$ of $S$, so every approximating factor is supported
on $N_{2r}(S)$.  This completes the proof.
\end{proof}

Ground-state-preserving products that suppress excited states also occur
in detectability-lemma constructions~\cite[Section~III]{AnshuAradVidick2016}.
For the random contractions here, the gap gives the following estimate.  Under
\eqref{eq:amplification-positive-system}, for any vector $\zeta$
including spectator systems,
\begin{equation}\label{eq:amplification-excited-decay}
 \mathbb E\|D(I-P_\Omega)\zeta\|^2
 \le e^{-gt}\|(I-P_\Omega)\zeta\|^2.
\end{equation}
Here and below physical projections are tensored with the identity
on spectators.  To see this, every self-adjoint $G_i$ fixes $\Omega$
and commutes with $P_\Omega$.  At a clock event the squared norm of
an excited vector decreases by $\langle\zeta,k_i\zeta\rangle$.
Summing the conditional rate of these decreases gives at least
$g\|\zeta\|^2$, and integrating proves
\eqref{eq:amplification-excited-decay}.  Thus the global gap is
used for decay, while the preceding channel argument supplies locality.

\subsection{A local unitary from a small collar cost}

For an ambient set $T\subset\Z^2$, write
$T_L=T+[-L,L]^2_{\Z}$ as in the scanner.  The next proposition
states explicitly the information it needs from that scanner.

\begin{proposition}[Amplification]\label{prop:amplification}
Fix $C_0>0$ and let $\varepsilon,\beta,\alpha$ be as in
\eqref{eq:amplification-exponents}.  Suppose $n\ge2$,
$L=\lfloor n^{1-\varepsilon}\rfloor$, and
\[
 X=A\cap T,\qquad S_0=A\cap T_L,\qquad Y=S_0\setminus X,
 \qquad Z'=\Lambda\setminus S_0.
\]
Assume that $|S_0|\le C_0n^2$, that each of the three cuts
$X,Y,S_0$ has at most $C_0n$ crossing edges, and that
\begin{equation}\label{eq:amplification-input-information}
 \mathcal I:=I_\Omega(X:Z')\le C_0n^{1-\varepsilon}.
\end{equation}
For every integer $M\ge1$, at all sufficiently large $n$, there is
an integer radius $r$ such that
\begin{equation}\label{eq:amplification-radius}
 r\le C_M n^{(1-\varepsilon)/\alpha},\qquad
 L+2r=o(n^\beta),
\end{equation}
and, simultaneously for every physical set
$J\subseteq\Lambda\setminus N_{2r}(S_0)$,
\begin{equation}\label{eq:amplification-output}
 I_\Omega(X:J)\le C_M n^{-M}.
\end{equation}
The same bound holds with any subset of $X$ in place of $X$.
The constants and the threshold may depend on $C_0$ and $M$, but are
independent of $\Lambda,A,T$ subject to the stated assumptions.

If in addition
\begin{equation}\label{eq:amplification-cut-clearance}
 \dist_\infty(T,Z)>L+n^\beta,
\end{equation}
where $Z$ is the set of endpoints of the original cut, then
$N_{2r}(S_0)\subseteq A$.  In particular the remote set $J$ may
include all of $\Lambda\setminus A$.
\end{proposition}

\begin{proof}
The case $X=\varnothing$ is immediate; take $r=1$.  Suppose that
$X$ is nonempty.  We keep the entire physical system fixed while
constructing the auxiliary operators; the estimates below are uniform
in its volume.

Set $w=n^{3/5}$.  Lemma~\ref{lem:tail}, applied in the original
ground state to the three cuts, provides spectral projections
$P_X,T_Y,Q_{S_0}$ onto their positive eigenvalues with surprisal in
the intervals
\[
 [S(X)-w,S(X)+w],\quad [S(Y)-w,S(Y)+w],\quad
 [S(S_0)-w,S(S_0)+w],
\]
respectively.  The original interactions have uniformly bounded
support dimensions and at most $Cn$ terms cross each cut.
Thus each failed projection has probability at most
$C\exp(-cn^{1/10})$.  For any fixed integer $K\ge1$, increasing
the scale threshold makes these three probabilities at most
\begin{equation}\label{eq:amplification-typical-error}
 \delta=n^{-2K}.
\end{equation}
For a trivial subsystem its projection is the identity.  All inverse
eigenvalues used below are restricted to the selected positive
spectrum, so no invertibility assumption is imposed on a marginal.

Choose a full eigenbasis $\{|x\rangle\}$ for $\rho_X$, with
eigenvalues $p_x\ge0$, and write a Schmidt expansion
\[
 \Omega=\sum_{y:\,p_y>0}\sqrt{p_y}\,
 |y\rangle_X|\omega_y\rangle_{YZ'}.
\]
Append two systems $s',e'$, each a copy of the Hilbert space on $X$,
with fixed unit blank vectors denoted $|0\rangle$.
Consider the unit vectors
\begin{align}
 \phi&=|0\rangle_X|0\rangle_{s'}
       \sum_{y:\,p_y>0}\sqrt{p_y}|y\rangle_{e'}
                                    |\omega_y\rangle_{YZ'},
       \label{eq:amplification-swapped-state}\\
 \eta&=\Omega\otimes
       \sum_{x:\,p_x>0}\sqrt{p_x}|x\rangle_{s'}|x\rangle_{e'}.
       \label{eq:amplification-restored-state}
\end{align}
In $\phi$, the original target has moved to $e'$, so its joint
reduction with a remote physical system is unchanged after relabeling.
In $\eta$, $e'$ is purified by $s'$ and is independent of all physical
systems.  The unitary we construct will leave $e'$ untouched and act
only on $s'$ and a graph neighborhood of $S_0$.

For readability write $Q=Q_{S_0}$ and $T_0=T_Y$.  On $s',X,Y$
define the square operator
\begin{equation}\label{eq:amplification-restoring-operator}
 R_0=\sum_{x\ {\rm selected}}p_x^{-1/2}
 |x\rangle\langle0|_{s'}\otimes
 Q\bigl(|x\rangle\langle0|_X\otimes T_0\bigr).
\end{equation}
It acts trivially on $e'$ and $Z'$.  The localization lemma applies to
vectors of the form $B(\Omega\otimes\xi)$ with a local operator $B$.
To put $R_0\phi$ in this form, define an operator supported on
$S_0,s',e'$ by
\begin{equation}\label{eq:amplification-column-operator}
 B_0=\sum_{\substack{x\ {\rm selected}\\y\ {\rm in\ the\ full\ basis}}}
 p_x^{-1/2}|x\rangle\langle0|_{s'}\otimes
 |y\rangle\langle0|_{e'}\otimes
 Q\bigl(|x\rangle\langle y|_X\otimes T_0\bigr).
\end{equation}
The summation over $y$ includes zero-probability eigenvectors.
Direct substitution gives
\begin{equation}\label{eq:amplification-column-identity}
 R_0\phi=B_0(\Omega\otimes|00\rangle_{s'e'}).
\end{equation}

We establish three estimates for these operators before introducing
the constraint product.  The first is a norm bound controlled by
$\mathcal I$.  Let
\[
 B_Y=\sum_{x\ {\rm selected}}p_x^{-1}
       T_0\langle x|Q|x\rangle T_0\ge0.
\]
Orthogonality of the output vectors in
\eqref{eq:amplification-restoring-operator} gives
$R_0^\dagger R_0=|0\rangle\langle0|_{s'}\otimes
|0\rangle\langle0|_X\otimes B_Y$.
Similarly,
$B_0^\dagger B_0=|00\rangle\langle00|_{s'e'}\otimes
I_X\otimes B_Y$.
The typical spectral inequalities are
\[
 p_x^{-1}\le e^{S(X)+w},\qquad
 Q\le e^{S(S_0)+w}\rho_{S_0},\qquad
 T_0\rho_Y T_0\le e^{-S(Y)+w}T_0.
\]
They imply
\begin{equation}\label{eq:amplification-operator-norm}
 \|R_0\|,\ \|B_0\|\le M_0,\qquad
 M_0=\exp\bigl((\mathcal I+3w)/2\bigr),
\end{equation}
because purity gives
$S(X)+S(S_0)-S(Y)=I(X:Z')=\mathcal I$.

The second estimate is an adjoint identity with no factor $M_0$.
Let $\mathcal R_X$ be the isometry that relabels $X$ as $e'$ and
inserts the blank vectors on $X,s'$; thus $\mathcal R_X\Omega=\phi$.
For every physical operator $L'$ the coefficients $\sqrt{p_x}$
in \eqref{eq:amplification-restored-state} cancel the selected
inverse square roots, giving exactly
\begin{equation}\label{eq:amplification-adjoint-identity}
 R_0^\dagger (L')^\dagger\eta
 =\mathcal R_X(P_X\otimes T_0)Q(L')^\dagger\Omega.
\end{equation}
The three projections are contractions, whether or not they commute.
Telescoping them on $\Omega$ and using
\eqref{eq:amplification-typical-error} therefore gives
\begin{equation}\label{eq:amplification-adjoint-error}
 \|R_0^\dagger(L')^\dagger\eta-\phi\|
 \le3\sqrt\delta+\|(L')^\dagger\Omega-\Omega\|.
\end{equation}

Third, although $R_0\phi$ may have large norm, its physical ground
component has norm at most one.  Define
\[
 M_X=\Tr_Y\bigl[(I_X\otimes T_0)\rho_{S_0}Q\bigr].
\]
Since $Q$ commutes with $\rho_{S_0}$, one has
$0\le\rho_{S_0}Q\le\rho_{S_0}$.
Cyclicity of the partial trace for operators on $Y$ allows a second
$T_0$ on the right.  Hence
\[
 0\le M_X\le\Tr_Y[(I\otimes T_0)\rho_{S_0}(I\otimes T_0)]
 \le\rho_X.
\]
The last inequality follows by adding the analogous positive term
with $I-T_0$; the two cross terms have zero partial trace.
In particular $M_X$ vanishes on $\ker\rho_X$.  On its support put
$A_X=\rho_X^{-1/2}M_X\rho_X^{-1/2}$, so $0\le A_X\le I$.
With $\rho_X^+$ denoting the inverse on the support and zero on the
kernel,
\[
 M_X\rho_X^+M_X
 =\rho_X^{1/2}A_X^2\rho_X^{1/2}\le M_X.
\]
The ancillary coefficient of $|x\rangle_{s'}|y\rangle_{e'}$
in $P_\Omega R_0\phi$ is $(M_X)_{yx}/\sqrt{p_x}$ for selected
$x$.  It follows that
\begin{equation}\label{eq:amplification-ground-component}
 \|P_\Omega R_0\phi\|^2
 =\sum_{x\ {\rm selected},y}\frac{|(M_X)_{yx}|^2}{p_x}
 \le\Tr(M_X\rho_X^+M_X)\le\Tr M_X\le1.
\end{equation}

We have now constructed a restoring operator whose norm is
$e^{O(\mathcal I+w)}$, whose adjoint is accurate on the desired
state, and whose ground component is bounded.  Positive constraints
will suppress the remaining component while preserving the last two
properties.

Run the clocks of Lemma~\ref{lem:amplification-local-products} for
time
\begin{equation}\label{eq:amplification-time-choice}
 t_*=C(\mathcal I+w+K\log n),\qquad
 r=\left\lceil\bigl(C'(t_*+K\log n)\bigr)^{1/\alpha}\right\rceil,
\end{equation}
where $C$ and then $C'$ are chosen sufficiently large.  Use $S=S_0$,
and put $\zeta=R_0\phi$.  The full, retained, and truncated products
are denoted $D,D_{\rm ret},L'$ respectively.
Equations~\eqref{eq:amplification-excited-decay} and
\eqref{eq:amplification-operator-norm} imply
\[
 \mathbb E\|D(I-P_\Omega)\zeta\|^2
 \le M_0^2e^{-gt_*}.
\]
Apply Lemma~\ref{lem:amplification-local-products} with
$B=B_0$, $\xi=|00\rangle_{s'e'}$, and $S=S_0$.
The column representation \eqref{eq:amplification-column-identity}
then gives
\begin{align*}
 \mathbb E\|(D-D_{\rm ret})\zeta\|
 &\le Ct_*|S_0|M_0e^{vt_*-cr^\alpha},\\
 M_0\mathbb E\|D_{\rm ret}-L'\|
 &\le Ct_*|S_0|(1+r)^2M_0e^{-cr^\alpha}.
\end{align*}
All of these estimates are on the same clock space.

Here is explicit control of the scale choices.  First choose $C$
so that $gt_*\ge\mathcal I+3w+6K\log n$.
The first expectation is then at most $n^{-6K}$.
By \eqref{eq:amplification-input-information}, both $t_*$ and $r$
are bounded by fixed powers of $n$, with constants depending on $K$.
Their polynomial prefactors and $|S_0|$ therefore cost only
$O_K(\log n)$ in the logarithm.  Since
$\log M_0=(\mathcal I+3w)/2\le Ct_*$, choosing $C'$ larger
makes each of the last two expectations at most $n^{-3K}$ for
large $n$.  Markov's inequality and a union bound consequently
provide a realization satisfying
\begin{equation}\label{eq:amplification-good-trajectory}
 \|D(I-P_\Omega)\zeta\|\le n^{-K},\quad
 \|(D-D_{\rm ret})\zeta\|\le n^{-K},\quad
 M_0\|D_{\rm ret}-L'\|\le n^{-K}.
\end{equation}
Indeed the three failure probabilities sum to at most
$n^{-4K}+2n^{-2K}<1$ at sufficiently large scales.

Every exact factor fixes the ground vector in both directions.
Thus $D$ preserves the ground component in
\eqref{eq:amplification-ground-component}, and
$D_{\rm ret}^\dagger\Omega=\Omega$.
Since $M_0\ge1$, \eqref{eq:amplification-good-trajectory} gives
\begin{equation}\label{eq:amplification-before-polar}
 \|L'R_0\phi\|\le1+3n^{-K},\qquad
 \|(L')^\dagger\Omega-\Omega\|\le n^{-K}.
\end{equation}
Combining the latter bound with
\eqref{eq:amplification-adjoint-error}, and writing $R'=L'R_0$,
we obtain
\[
 \|(R')^\dagger\eta-\phi\|\le4n^{-K}.
\]
Because $\phi,\eta$ are unit vectors, the real overlap satisfies
$\Re\langle\eta,R'\phi\rangle\ge1-4n^{-K}$.
Expanding the square using
\eqref{eq:amplification-before-polar} gives
\begin{equation}\label{eq:amplification-direct-error}
 \|R'\phi-\eta\|\le Cn^{-K/2}.
\end{equation}

The square operator $R'$ is supported on $s'$ and $N_{2r}(S_0)$,
and acts trivially on $e'$.  Its polar decomposition extends on this
finite local space to
$R'=U'H_+$, where $U'$ is unitary and $H_+=(R'^\dagger R')^{1/2}$.
Such a unitary extension exists even when $R'$ is singular: its two
polar support spaces have equal dimension, as do their orthogonal
complements.  Tensor it with the identity elsewhere.
For $\xi=(U')^\dagger\eta$ we have
\[
 \|H_+\xi-\phi\|\le4n^{-K},\qquad
 \|H_+\phi-\xi\|\le Cn^{-K/2}.
\]
Subtract these two errors in the identity
\[
 (I+H_+)(\xi-\phi)
 =(H_+\xi-\phi)-(H_+\phi-\xi).
\]
Since $\|(I+H_+)^{-1}\|\le1$, this proves
\begin{equation}\label{eq:amplification-unitary-error}
 \|U'\phi-\eta\|\le Cn^{-K/2}.
\end{equation}
In particular no bound on $\|H_+\|$ is needed.

Let $J\subseteq\Lambda\setminus N_{2r}(S_0)$ be arbitrary.
The unitary $U'$ changes neither $e'$ nor $J$, so their joint
reduction in $U'\phi$ is exactly the original $XJ$ reduction,
with $X$ relabeled as $e'$.  In $\eta$ that joint reduction is a
product.  Apply Lemma~\ref{lem:continuity} to
$I(e':J)=S(e')-S(e'|J)$ and
\eqref{eq:amplification-unitary-error}.  Only the dimension of
$e'$ appears.  It has
$\log\dim e'=|X|\log q\le Cn^2$, and therefore
\[
 I_\Omega(X:J)\le Cn^{2-K/4}.
\]
Choose the fixed integer $K\ge4(M+3)$.
This proves \eqref{eq:amplification-output}, uniformly for every
such $J$.  Discarding part of $X$ proves the assertion for its subsets.

Finally \eqref{eq:amplification-time-choice} and
\eqref{eq:amplification-input-information} give
$r\le C_M n^{(1-\varepsilon)/\alpha}$.
The choices in \eqref{eq:amplification-exponents} imply
$(1-\varepsilon)/\alpha<\beta$ and $1-\varepsilon<\beta$,
which prove \eqref{eq:amplification-radius}.
If a graph path of length at most $2r$ from $S_0$ reached
$\Lambda\setminus A$, it would contain an endpoint of the original
cut.  A starting point of this path has ambient distance at most
$L$ from $T$, and each graph step changes ambient distance by at
most one.  That endpoint would be within $L+2r$ of $T$,
contradicting \eqref{eq:amplification-cut-clearance} for large $n$.
Thus the graph neighborhood lies in $A$, as claimed.
\end{proof}

The global gap enters through the marginal-tail estimate, the
construction of positive constraints, and the excited-state decay.
Its final entropy comparison is confined to the target dimension.
The next section constructs tiles whose enlarged graph neighborhoods
avoid all previously used tiles of the same family, so that
Proposition~\ref{prop:amplification} can be applied repeatedly.
